\chapter{Machine Learning Methodology}
\label{chap:ml_methodology}

\par This chapter presents the machine learning core of the thesis: how the raw CICIoT2023 captures are turned into clean training data, how that data is split so that the evaluation is trustworthy, how the model is built and trained, and how it is evaluated. Everything else in the system is built on top of the artefacts this chapter produces. The application that serves the trained model is specified in Chapter~\ref{chap:requirements} and designed in Chapter~\ref{chap:design_implementation}; the experimental results of the methodology described here are reported in Chapter~\ref{chap:results}.

\section{Classification Tasks}
\label{sec:classification_tasks}

\par The dataset provides 34 fine-grained labels (the 33 attack types of Section~\ref{subsec:attack_taxonomy} plus benign traffic). This thesis does not classify at that finest level. Instead, the same rows are relabelled into two classification tasks of practical interest:

\begin{itemize}
    \item \textbf{Binary (2-class)}: benign versus attack. This is the task an alerting IDS uses to decide whether traffic is malicious at all, and it is the primary task of the thesis.
    \item \textbf{8-class}: the seven attack families of Table~\ref{tab:attack_taxonomy} (DDoS, DoS, Reconnaissance, Web-based, Brute Force, Spoofing, Mirai) plus benign. This is the coarser attack-family view an analyst consumes, reported as a secondary task.
\end{itemize}

\par Both tasks are derived from the same subsampled rows and the same train/validation/test partitions; only the label mapping changes. The 34-class fine-grained task is not modelled, because the rarest variants hold too few unique records (Section~\ref{subsec:dataset_stats}) to train and evaluate reliably, and because attack-family granularity is sufficient for the alerting and triage use cases the demonstration targets.

\section{Data Pipeline}
\label{sec:data_pipeline_design}

\par The training data pipeline performs a sequence of deterministic transformations from the raw per-folder CSVs to the train/validation/test tensors that feed the model. Each transformation is a separate pipeline phase so that intermediate state is inspectable.

\subsection{Ingestion and Deduplication}
\label{subsec:ingestion}

\par For each of the attack folders, the pipeline lazily scans every CSV, projects only the feature columns plus a synthetic column that carries the source-capture filename, and concatenates the resulting lazy frames. Materialisation happens only after the projection, so the multi-tens-of-millions-of-rows corpus is never held in memory in its raw width.

\par CICIoT2023 contains a substantial fraction of exact duplicate rows, a known property of the dataset and a recurring concern in the literature. Deduplication is applied immediately after concatenation, before any sampling. Deduplicating before sampling is essential: a folder where most rows are duplicated would otherwise have the same row repeated many times in the sample, which would inflate accuracy on the corresponding test rows. The quantitative impact of each pipeline stage (raw count, duplicates removed, unique rows, sampling cap applied) is reported in Chapter~\ref{chap:results} (Figure~\ref{fig:ingest_waterfall}).

\subsection{Per-Class Subsampling}
\label{subsec:subsampling}

\par The \emph{per-class cap} is a fixed upper bound on the number of rows kept for each fine-grained class, set to 200{,}000 rows in this work. After deduplication, every folder whose deduplicated row count exceeds this cap is randomly subsampled down to it with a fixed seed, while folders below the cap are kept in full. This brings the effective training-set size down from tens of millions of rows to a few million while keeping every attack class represented.

\par The cap is set at the high end of the range used in the CICIoT2023 literature. Random sampling is preferred over taking the leading rows of each folder because each attack folder is a chronological concatenation of several capture sessions: the first $N$ rows come almost entirely from the first session, and a leading-rows sample would bias the model toward whatever happened in that opening session.

\subsection{Caching}
\label{subsec:caching}

\par The deduplicated, subsampled, labelled rows are concatenated across folders and written to a single Zstandard-compressed columnar file. Subsequent runs skip ingestion and load from this cache directly. Regeneration is forced by deleting the cache file. This is documented inline so future runs do not silently use a stale cache after sampling logic changes.

\subsection{Cleaning}
\label{subsec:cleaning}

\par The cleaning phase performs three operations on the materialised frame. First, null-label rows are dropped (defensive only; labels come from folder names, so this should be a no-op). Second, infinities are replaced with nulls; the actual median imputation is deferred to after the train/validation/test split so the medians are computed on the training partition only. Third, a $\log(1+x)$ transform is applied to the heavy-tailed continuous features (rates, total bytes, inter-arrival times). Binary protocol and flag indicators are excluded from the log transform.

\subsection{Data Split}
\label{subsec:split_strategies}

\par All experiments in this thesis use a single split type, the stratified random split. The subsampled rows are partitioned into 70/15/15 train/validation/test sets by two stratified splits in series, stratified by the fine-grained label so that every one of the 34 classes appears in all three partitions. This random split follows the convention used by most published CICIoT2023 studies, which keeps the reported numbers comparable to that body of work.

\par A random row split is known to be optimistic. Because each attack folder is a contiguous capture of a single session, two flows recorded moments apart can land on opposite sides of the train/test boundary, so some test rows have near-duplicate neighbours in training and the measured performance is an upper bound on what the model would achieve on genuinely unseen traffic. Two leakage-aware alternatives, a group-wise split that keeps each capture whole and a temporal split that trains on earlier captures and tests on later ones, would give a more conservative estimate; neither is used here, and both are identified as future work in Chapter~\ref{chap:conclusions}.

\subsection{Train-Only Preprocessing Fit}
\label{subsec:train_fit_preprocess}

\par For the split, column medians are computed on the training rows only, nulls in the train, validation, and test partitions are imputed using those medians, and the robust scaler is fitted on training rows alone before being applied to all three partitions. The robust scaler standardises each feature by its median and interquartile range rather than its mean and standard deviation:
\begin{equation}
    x' = \frac{x - \operatorname{median}(x)}{\mathrm{IQR}(x)}, \qquad \mathrm{IQR}(x) = Q_3(x) - Q_1(x),
    \label{eq:robust_scaler}
\end{equation}
where $Q_1$ and $Q_3$ are the first and third quartiles estimated on the training rows. Centring on the median and dividing by the IQR makes the transform insensitive to the heavy upper tails of the CIC rate and byte-count features, which would otherwise dominate a mean/variance scaler. The contract that medians, quartiles, and the scaler are fitted strictly on training rows rules out preprocessing leakage from the validation and test partitions into the fit.

\section{Model Architecture and Training}
\label{sec:model_architecture}

\subsection{MLP Topology}
\label{subsec:mlp_topology}

\par The model is a feed-forward multi-layer perceptron whose hidden layers each apply a linear transform followed by batch normalisation, a non-linear activation, and dropout, ending in a linear projection to the class scores. The number and width of the hidden layers, the dropout rate, and the activation are not fixed by hand: they are selected per task by the hyperparameter search of Section~\ref{subsec:hpo}, with the resulting configurations listed in Table~\ref{tab:hpo_selected}. Both selected networks take the 25-feature input vector and apply batch normalisation and dropout after every hidden layer:

\begin{align*}
    \text{binary (2-class):}\quad & 25 \to 512 \to 192 \to 512 \to 2 \\
    \text{8-class:}\quad & 25 \to 352 \to 256 \to 320 \to 64 \to 8
\end{align*}
where each arrow other than the final one denotes the Linear--BatchNorm--activation--Dropout block above.

\par Batch normalisation \cite{Ioffe2015}, applied after each hidden linear layer, standardises that layer's pre-activations across the mini-batch to zero mean and unit variance before rescaling them by two learned parameters $(\gamma, \beta)$:
\begin{equation}
    \mathrm{BN}(z) = \gamma \, \frac{z - \mu_{\mathcal{B}}}{\sqrt{\sigma_{\mathcal{B}}^2 + \epsilon}} + \beta,
    \label{eq:batchnorm}
\end{equation}
where $\mu_{\mathcal{B}}$ and $\sigma_{\mathcal{B}}^2$ are the per-feature mean and variance over the current batch $\mathcal{B}$ and $\epsilon$ is a small constant for numerical stability.

\par Two design elements are held fixed across both networks regardless of the searched depth and width. Batch normalisation (Equation~\ref{eq:batchnorm}) after every hidden layer stabilises training under the large class-weight range: at the 8-class granularity the heaviest class weight is several times the lightest, and gradients are noisier without it. Dropout after every hidden layer provides regularisation; the search settled on moderate rates (around 0.10--0.13, Table~\ref{tab:hpo_selected}) rather than the more aggressive 0.5 often used on smaller datasets, which is consistent with the cleaned training set already holding several million rows.

\subsection{Class-Weighted Cross-Entropy}
\label{subsec:class_weighted_ce}

\par Class weights are computed using the standard balanced-frequency rule applied to the encoded training labels, then passed to the cross-entropy loss. For a problem with $K$ classes, $N$ training samples, and $n_c$ samples in class $c$, the weight of class $c$ is
\begin{equation}
    w_c = \frac{N}{K \, n_c},
    \label{eq:balanced_weights}
\end{equation}
so that rare classes receive proportionally larger weight (and $\sum_c n_c w_c = N$). These weights enter the multi-class cross-entropy loss applied to the softmax outputs $\hat{y}_{i,k}$:
\begin{equation}
    \mathcal{L} = -\frac{1}{N} \sum_{i=1}^{N} \sum_{k=1}^{K} w_{k}\, y_{i,k} \log\!\big(\hat{y}_{i,k}\big),
    \label{eq:weighted_ce}
\end{equation}
where $y_{i,k}$ is the one-hot ground-truth indicator for sample $i$ and class $k$. Equation~\eqref{eq:weighted_ce} is the multi-class, class-weighted generalisation of the binary cross-entropy in Equation~\eqref{eq:bce_loss}. Weighting is preferred over resampling because SMOTE would generate synthetic flow vectors that have no operational meaning (a half-and-half SYN-flood / DNS-spoofing flow does not correspond to any real attack), and weighted batch sampling would be redundant with weighted cross-entropy.

\subsection{Optimiser, Schedule, and Early Stopping}
\label{subsec:training_loop}

\par Training uses the AdamW optimiser for both tasks, with the initial learning rate and the batch size selected per task by the search (Table~\ref{tab:hpo_selected}). A learning-rate schedule halves the rate after two consecutive epochs without validation-loss improvement, and early stopping with a patience of five epochs caps training at a maximum of fifty epochs. The batch sizes used (2048 for the binary task and 4096 for the 8-class task) are large enough to amortise the batch-normalisation overhead while keeping the per-step gradient meaningful.

\par The training loop tracks training loss, validation loss, and the current learning rate per epoch, restoring the best validation-loss state at the end. The best state is what is serialised; the final-epoch state is discarded.

\subsection{Hyperparameter Selection}
\label{subsec:hpo}

\par The architecture and optimisation settings above are not arbitrary: they were selected by an automated hyperparameter search using the Optuna framework \cite{Akiba2019} rather than fixed by hand. A Tree-structured Parzen Estimator (TPE) sampler explores the search space summarised in Table~\ref{tab:hpo_space} over thirty trials, each trained for ten epochs on a 100{,}000-row stratified subsample of the training set. The objective maximised is the validation macro-F1, chosen over validation loss so the search rewards balanced performance across the imbalanced classes rather than overall likelihood. The search is run separately for each task, so the selected configurations differ by task; they are listed in Table~\ref{tab:hpo_selected}. Each selected configuration is then retrained to convergence (up to fifty epochs with early stopping) on the full training set and used for all reported results. The Random Forest baseline is tuned by the same Optuna TPE procedure over thirty trials, also maximising validation macro-F1, and its selected hyperparameters appear alongside the MLP in Table~\ref{tab:hpo_selected}.

\begin{table}[htbp]
\begin{center}
\begin{tabular}{|p{120pt}|p{200pt}|}
\hline
\textbf{Hyperparameter} & \textbf{Search range} \\
\hline
\hline Hidden layers & 2--4 \\
\hline Units per layer & 64--512 (step 32) \\
\hline Dropout & 0.1--0.5 \\
\hline Activation & ReLU, ELU, LeakyReLU \\
\hline Learning rate & $10^{-4}$--$10^{-2}$ (log scale) \\
\hline Optimiser & Adam, AdamW \\
\hline Batch size & 2048, 4096, 8192 \\
\hline
\end{tabular}
\end{center}
\caption{Optuna TPE search space for the MLP. Thirty trials, ten epochs each on a 100{,}000-row subsample, validation macro-F1 as the objective.}
\label{tab:hpo_space}
\end{table}

\begin{table}[htbp]
\begin{center}
\begin{tabular}{|p{150pt}|p{90pt}|p{90pt}|}
\hline
\textbf{Hyperparameter} & \textbf{Binary (2-class)} & \textbf{8-class} \\
\hline
\multicolumn{3}{|l|}{\textbf{MLP}} \\
\hline
Hidden layer widths & 512, 192, 512 & 352, 256, 320, 64 \\
\hline Dropout & 0.13 & 0.10 \\
\hline Activation & ReLU & LeakyReLU \\
\hline Optimiser & AdamW & AdamW \\
\hline Learning rate & $1.6\times10^{-3}$ & $9.5\times10^{-3}$ \\
\hline Batch size & 2048 & 4096 \\
\hline
\multicolumn{3}{|l|}{\textbf{Random Forest}} \\
\hline
Trees ($n_\text{estimators}$) & 300 & 450 \\
\hline Max depth & 25 & 25 \\
\hline Min samples split & 7 & 6 \\
\hline Min samples leaf & 3 & 4 \\
\hline Max features & 0.7 & 0.7 \\
\hline Class weight & balanced & balanced \\
\hline
\end{tabular}
\end{center}
\caption{Hyperparameters selected by the Optuna TPE search for each task, used for all reported results. The Random Forest \texttt{max\_depth} is capped at 25 to keep the serialised forest tractable.}
\label{tab:hpo_selected}
\end{table}

\section{Evaluation Methodology}
\label{sec:evaluation_methodology}

\subsection{Classification Metrics}
\label{subsec:classification_metrics}

\par Every (granularity, model) combination is evaluated under the same protocol: accuracy, precision, recall, and F1, reported both per class and as macro and support-weighted averages, alongside the row-normalised confusion matrix. The four base metrics derive from the per-class one-versus-rest counts: true positives (TP), false positives (FP), true negatives (TN), and false negatives (FN). Accuracy is the fraction of flows classified correctly; precision is the share of flows predicted as a class that truly belong to it; recall is the share of a class's flows that are recovered; and F1 is the harmonic mean of precision and recall, which stays low unless both are high:
\begin{equation}
\begin{aligned}
\mathrm{Accuracy}  &= \frac{TP + TN}{TP + TN + FP + FN}, &\qquad
\mathrm{Precision} &= \frac{TP}{TP + FP}, \\[8pt]
\mathrm{Recall}    &= \frac{TP}{TP + FN}, &\qquad
\mathrm{F1}        &= \frac{2\,\mathrm{Precision}\cdot\mathrm{Recall}}{\mathrm{Precision} + \mathrm{Recall}}.
\end{aligned}
\end{equation}

\par For the multi-class tasks each score is averaged over classes two ways, because they answer different questions under class imbalance. The \emph{weighted} average weights each class's score by its support, so it reflects the traffic mix as it actually occurs but can be dominated by the large DDoS and Mirai classes. The \emph{macro} average weights every class equally, so it exposes weakness on the rare classes that the weighted average can hide. Aggregate numbers are always accompanied by the per-class report.

\subsection{Confidence Calibration}
\label{subsec:calibration_method}

\par Each verdict is shown with a confidence score, the largest of the model's softmax probabilities. For that number to be useful it should mean what it says: among the flows a model calls attacks with 90\% confidence, about 90\% should truly be attacks. Modern neural networks rarely satisfy this out of the box; they tend to be \emph{overconfident}, reporting peak probabilities well above their actual accuracy \cite{Guo2017}. The reason lies in the training objective. Cross-entropy is minimised by driving the correct-class probability towards one, and once a flow is already classified correctly the only way to reduce its loss further is to widen the gap between the correct logit and the rest. Training therefore keeps inflating logit magnitudes long after the predictions have stopped changing, and the softmax peak saturates towards one. Class-weighted training sharpens this on the classes it emphasises, which is why (Section~\ref{sec:calibration}) the binary model, dominated by a single heavily-weighted attack class, ends up markedly more overconfident than the eight-class model.

\par \emph{Temperature scaling} corrects the scores without touching the predictions. A single scalar $T > 0$ is fitted on the validation-set logits $\mathbf{z}$ by minimising the negative log-likelihood, and the served probabilities are recomputed as $\mathrm{softmax}(\mathbf{z}/T)$; a value $T > 1$ flattens the distribution and pulls the inflated peak back down towards the empirical accuracy. The property that makes this safe is monotonicity: dividing every logit by the same positive constant leaves their ordering intact, so the largest logit stays the largest. The predicted label is therefore unchanged, and with it the accuracy, precision, recall, and the entire confusion matrix; only the reported probability moves. Calibration is in this sense a free correction, improving the honesty of the confidence number at no cost to any classification metric. It also explains why a single parameter is preferred over the richer per-class scalings, which have enough freedom to overfit the validation set and perturb the predictions \cite{Guo2017}.

\par Calibration quality is measured by the Expected Calibration Error (ECE): predictions are sorted into fifteen equal-width confidence bins, and ECE is the support-weighted average gap between the mean confidence and the mean accuracy within each bin, so a perfectly calibrated model scores zero. The fitted temperatures and the resulting change in ECE are reported in Section~\ref{sec:calibration}. The distinction is not cosmetic for an alerting tool: an operator who triages by confidence, or suppresses alerts below a threshold, is trusting that a high score really does mean a high chance of attack, and an uncalibrated 99\% that is right only four times in five would quietly undermine that decision.

\subsection{Latency Benchmark}
\label{subsec:latency_benchmark}

\par For each trained model, the benchmark loop runs $1{,}000$ inference passes per batch size in $\{1, 32, 256, 1024\}$ on CPU after a warmup of one hundred passes, with the timer wrapped around the feature scaling step \emph{and} the model forward pass. The pipeline reports the p50, p95, and p99 of the per-batch latency in milliseconds and the implied flows-per-second throughput. The acceptance criterion, later formalised as NFR-1 (Section~\ref{sec:non_functional_requirements}), is a sub-100~ms p95 end-to-end latency on a consumer laptop CPU at any of those batch sizes.
